\section{模型的检验}

\subsection{需求预测模型检验}

需求预测按时间划分训练期与留出期。2019年3月1日至2024年3月31日用于估计日历趋势、星期效应与模型参数，2024年4月1日至2025年3月31日仅用于评价。滚动8周模型在每个预测日只读取此前数据。采用WAPE和sMAPE比较7日季节朴素、过去8周同星期中位数与日历趋势对数岭回归，结果见表\ref{tab:forecast-comparison}。住院和门诊由8周同星期中位数取得最低WAPE，分别为16.04\%和14.04\%；体检受团检批次影响，最低WAPE仍为69.83\%。因此，住院和门诊可由历史星期节律滚动预测，体检配置需要未来批次计划或高分位压力情景。该现象与医院门诊量中常见的趋势和周内效应一致\cite{zhu2015,xiang2009}。

模型参数的时间边界另行核对：峰平阈值、项目目录、医生容量和设备稀缺权重只从训练期提取；留出期不参与参数选择。报告间隔仅作项目族相对负荷代理，排程时长仍使用题给区间。由此，表\ref{tab:forecast-comparison}中的误差反映向后续日期外推的表现，而不是样本内拟合程度。

\subsection{调度模型有效性检验}

为检验调度对象的构造口径，重新核对事件键、项目数与设备集合。精确重复明细只计一次，同一患者同一开单时刻的不同项目仍归入同一事件；源系统显式连接的项目均按边界拆分。设备集合逐项对应题目表1，无法确认的项目保持空集合；床旁任务按式\eqref{eq:bedside-capability}取项目能力与7号机地点能力的交集。

问题二在患者顺序、候选项目序列、设备破同序、准备条件和转运规则均一致的条件下构造启发式对照。为使精确模型能够在有限时间内求解，本文从留出期数据按负荷、项目结构和设备约束抽取并额外构造6个代表性子问题；这些子问题是算法验证样本，并非赛题规定的运行情景。CP-SAT在每个子问题中最大化完整完成事件数，表\ref{tab:exact}同时报告精确完成数、最好界和求解状态；只有求解器状态为OPTIMAL时，精确完成数才作为该子问题的已证最优值。

\begin{table}[H]
\centering
\caption{问题二构造算法与CP-SAT精确子问题对照}
\label{tab:exact}
\begin{tabular}{lrrrrl}
\toprule
窗口 & 事件数 & 构造解 & CP-SAT & 最好界 & 状态 \\
\midrule
普通工作日 & 18 & 15 & 15 & 15 & OPTIMAL \\
高峰工作日 & 18 & 16 & 16 & 16 & OPTIMAL \\
周末 & 18 & 16 & 16 & 16 & OPTIMAL \\
多项目高峰 & 10 & 10 & 10 & 10 & OPTIMAL \\
床旁高峰 & 18 & 2 & 2 & 2 & OPTIMAL \\
产科高峰 & 18 & 16 & 16 & 16 & OPTIMAL \\
\bottomrule
\end{tabular}

\end{table}

6/6个子问题包含20--100个事件，均由求解器证明为OPTIMAL；构造算法的完整完成数与精确模型可行值逐个相同，最大保守至界差距为0.00\%。这说明本文构造算法在所抽取的普通、高峰、周末、多项目、床旁和产科局部样本中达到已证最优，但该结论不外推为全年大规模问题的最优性证明。

问题三的精确子问题同时包含住院、门诊、体检、背景最低完成数以及全部资源和准备约束。本文从留出期按需求结构和资源竞争程度抽取并额外构造8类代表性子问题，覆盖普通联合需求、背景高峰、住院高峰、多项目、床旁、专项设备和高竞争样本；它们用于验证算法，不是赛题预设情景。每个子问题按自身背景事件数施加与构造算法相同的局部$\alpha=0.45$约束，仅用于同规模解法对照；该缩小模型既不替代全年全局约束，也不表示正式模型设置了逐日配额。

\begin{table}[H]
\centering
\caption{问题三联合构造算法与CP-SAT精确子问题对照}
\label{tab:p3-exact}
\begingroup
\setlength{\tabcolsep}{3pt}
\begin{tabularx}{\textwidth}{>{\raggedright\arraybackslash}Xrrrrllrr}
\toprule
验证子问题 & $\alpha$ & 事件数 & \makecell{背景\\H/E} & \makecell{H\\住院} & \makecell{E住院值\\最好界} & 状态 & \makecell{求解器\\间隙} & \makecell{H至界\\差距} \\
\midrule
普通联合 & 0.45 & 20 & 12/7 & 4 & 4/4 & OPTIMAL & 0.00\% & 0.00\% \\
背景高峰 & 0.45 & 40 & 30/14 & 6 & 6/6 & OPTIMAL & 0.00\% & 0.00\% \\
住院高峰 & 0.45 & 60 & 24/11 & 32 & 32/32 & OPTIMAL & 0.00\% & 0.00\% \\
多项目高峰 & 0.45 & 40 & 19/11 & 14 & 14/14 & OPTIMAL & 0.00\% & 0.00\% \\
床旁压力 & 0.45 & 20 & 8/6 & 0 & 0/0 & OPTIMAL & 0.00\% & 0.00\% \\
产科压力 & 0.45 & 40 & 23/11 & 15 & 15/15 & OPTIMAL & 0.00\% & 0.00\% \\
设备竞争 & 0.45 & 60 & 34/17 & 23 & 23/23 & OPTIMAL & 0.00\% & 0.00\% \\
推荐点典型窗 & 0.45 & 100 & 60/27 & 33 & 33/33 & OPTIMAL & 0.00\% & 0.00\% \\
\bottomrule
\end{tabularx}
\endgroup

\end{table}

表\ref{tab:p3-exact}中H、E分别表示构造算法与精确模型，背景列按H/E并列；“E住院值/最好界”给出时限内精确模型的住院可行值$E$与最好界$U$。两类间隙均针对住院完成数目标：
\begin{align}
G_{\mathrm{solver}}&=\frac{U-E}{\max\{|U|,1\}},\\
G_{\mathrm{H,bd}}&=\frac{U-H}{\max\{|U|,1\}}.
\end{align}
当状态为OPTIMAL时，$E=U$，$G_{\mathrm{H,bd}}$就是构造解相对已证最优值的差距；当状态为FEASIBLE时，只报告可行值、最好界和两类保守间隙，不将可行值或最好界表述为最优解。8/8个代表性子问题均为OPTIMAL，构造解的住院完成数逐个等于已证最优值，最大保守至界差距为0.00\%。但这8个子问题的住院优先基线都已满足局部背景下限，没有触发全局缺口重构，因此上述结果只为基线分支提供局部最优证据，不能证明修复分支或全年联合问题达到全局最优。

为补充修复分支的实际路径证据，另按留出期题给上界工时选择2025年2月17日作为真实峰值日。该日共805个事件，其中住院305个、门诊407个、体检93个，上界工时合计20\,375分钟；局部背景事件为500个，$\alpha=0.45$对应最低完成数225。住院优先基线仅完成背景213个、住院259个，背景缺口为12，因而自然触发全局重构；稀缺能力优先修复完成背景258个、住院259个并被选中，14项独立检查的违规数均为0。该审计只证明真实数据中修复分支能够被自然触发并产生可行排程，不是修复分支的精确最优性证明。

为避免求解与验证共享同一段可行性判断，另从排程结果重建任务时长、释放与截止、设备能力、准备条件、设备互斥、患者互斥、跨室转运和医生并发约束。结果见表\ref{tab:verification}。问题二117\,843条排入任务与问题三159\,580条排入任务的各组违规数均为0，观察到的最大医生并发为20，未超过主情景容量。

\begin{table}[H]
\centering
\caption{问题二、三导出排程的独立可行性验证}
\label{tab:verification}
\begin{tabularx}{\textwidth}{X r l}
\toprule
核验项目 & 违规数 & 结论 \\
\midrule
任务键、输入对应、时长 & 0 & 通过 \\
设备兼容、床旁位置 & 0 & 通过 \\
释放、截止、工作时段 & 0 & 通过 \\
房间、患者、转运冲突 & 0 & 通过 \\
医生并发容量 & 0 & 通过 \\
结果表聚合一致性 & 0 & 通过 \\
\bottomrule
\end{tabularx}

\end{table}

现实校准按C0--C7逐层加入标准工时、工作时段、医生并发、项目设备能力、项目级准备和三类患者竞争，见表\ref{tab:calibration}与图\ref{fig:calibration}。历史开单至报告率为83.77\%，其时间语义与排程完成率不同，仅作为现实参照。仅加入题给工时后完成率为100\%，较历史口径高16.23个百分点；标准班次和训练期医生并发加入后仍为100\%。项目设备能力使完成率降至85.51\%，相邻下降14.49个百分点；项目级准备进一步使其降至85.30\%，相邻下降0.20个百分点。门诊、体检按$\alpha=0.45$进入共享日历后，住院率降至64.80\%，相邻下降20.50个百分点，说明专项设备集合和三类需求竞争是两项主要约束。C7仍为64.80\%、相邻变化0.00个百分点，因为该层只记录正式构造方法，没有再次改变可行域。

\begin{table}[H]
\centering
\caption{历史现实校准与约束增量}
\label{tab:calibration}
\begin{tabular}{lrrr}
\toprule
校准层 & 完成数/分母 & 完成率 & 较上一层/百分点 \\
\midrule
历史开单至报告 & 59,266/70,864 & 83.63\% & -- \\
仅题给时长 & 70,771/70,771 & 100.00\% & +16.37 \\
加入标准班次 & 70,771/70,771 & 100.00\% & +0.00 \\
加入医生并发 & 70,771/70,771 & 100.00\% & +0.00 \\
加入项目设备能力 & 60,526/70,771 & 85.52\% & -14.48 \\
加入项目级准备 & 60,359/70,771 & 85.29\% & -0.24 \\
加入门诊/体检竞争 & 42,730/70,771 & 60.38\% & -24.91 \\
正式构造策略 & 42,730/70,771 & 60.38\% & +0.00 \\
\bottomrule
\end{tabular}

\end{table}

\begin{figure}[H]
\centering
\includegraphics[width=\textwidth]{fig06_calibration.pdf}
\caption{从历史口径到正式联合排程的逐层现实校准}
\label{fig:calibration}
\end{figure}

\subsection{灵敏度分析}

保持评价期和推荐下限$\alpha=0.45$不变，分别改变医生容量、设备匹配边界、准备粒度、憋尿时间和转运时间，结果见表\ref{tab:sensitivity}和图\ref{fig:sensitivity}。每个情景只改变一个主因素，以便区分参数影响。表中同时给出$\varepsilon$约束是否满足及背景缺口；只有达到背景下限的情景才可作为同一政策约束下的可行方案比较，未达标情景只反映当前构造算法在该参数下的搜索结果。

\begin{table}[H]
\centering
\caption{推荐联合方案的单因素敏感性}
\label{tab:sensitivity}
\begin{tabular}{lrrrr}
\toprule
情景 & 住院48h率 & 背景预约日率 & 条件等待P90/h & 排入分钟 \\
\midrule
主情景：5\%上界时长 & 63.95\% & 59.70\% & 40.42 & 1,819,800 \\
无特殊时长上浮 & 63.97\% & 59.70\% & 40.42 & 1,819,660 \\
医生容量下四分位 & 61.96\% & 58.95\% & 42.00 & 1,769,275 \\
仅严格设备兼容 & 71.57\% & 49.67\% & 45.33 & 1,664,965 \\
憋尿准备45分钟 & 63.89\% & 59.88\% & 40.25 & 1,821,400 \\
憋尿准备90分钟 & 64.12\% & 59.17\% & 40.17 & 1,815,485 \\
跨室转运5分钟 & 63.93\% & 59.71\% & 40.50 & 1,821,800 \\
跨室转运15分钟 & 63.99\% & 59.76\% & 40.33 & 1,819,020 \\
\bottomrule
\end{tabular}

\end{table}

\begin{figure}[H]
\centering
\includegraphics[width=.88\textwidth]{fig07_sensitivity.pdf}
\caption{主要参数变化相对主情景的完成率差异}
\label{fig:sensitivity}
\end{figure}

八个单因素情景均满足$\varepsilon$约束，达标数为8/8，背景缺口均为0。医生容量降至训练期下四分位时，住院率和背景率分别为63.33\%和54.38\%，等待P90增至26.67小时，表明人员在岗容量下降会同时压低两类服务。只保留直接与明确语义匹配时，住院率为65.26\%，背景率降至47.98\%；住院率略升并不表示资质收紧改善系统，而是候选设备集合变化后，两类需求的竞争组合发生了改变。

事件级准备的住院率为63.02\%，低于项目级准备的64.80\%，背景率二者接近，分别为56.30\%和56.37\%；其等待P90由21.83小时升至24.83小时，支持将准备条件落实到对应项目。憋尿45/90分钟与转运5/15分钟四个边界情景中，两类完成率相对主情景的最大变化为2.22个百分点，出现在憋尿45分钟的背景率；同时该情景等待P90达到29.67小时。因此这些参数没有破坏$\alpha=0.45$下的可行性，但可能改变构造分支和等待尾部，不能只凭总体完成率判断其影响。

\subsection{稳健性分析}

复杂病例按题设比例抽取，采用五个固定随机种子重复运行。五次运行均满足$\varepsilon$约束，达标数为5/5，背景缺口均为0。住院率样本标准差仅为0.0114个百分点，取值为64.7935\%--64.8189\%，说明住院完成率对复杂病例具体落点较稳定；背景率样本标准差为0.9495个百分点，取值为56.3209\%--58.4584\%。等待P90范围为21.83--29.42小时，其中一个种子切换到稀缺能力优先的背景重构分支，产生更高的等待尾部。由此可认为$\alpha=0.45$下的约束可行性具有多种子稳定性，但背景服务量和等待尾部仍对离散构造路径敏感。

\begin{table}[H]
\centering
\caption{复杂病例多随机种子稳健性}
\label{tab:seed-robustness}
\begin{tabular}{lrrrr}
\toprule
指标 & 均值 & 标准差 & 最小值 & 最大值 \\
\midrule
住院48小时率 & 64.8039\% & 0.0114个百分点 & 64.7935\% & 64.8189\% \\
背景服务率 & 56.7601\% & 0.9495个百分点 & 56.3209\% & 58.4584\% \\
等待P90/小时 & 23.38 & 3.37 & 21.83 & 29.42 \\
排入分钟 & 1,932,254 & 81,073 & 1,895,420 & 2,077,280 \\
\bottomrule
\end{tabular}

\end{table}

需求状态方面，问题一用训练期分位数识别平峰、高峰与压力状态，留出期超过训练高峰阈值的日期占13.15\%，超过压力阈值的日期占7.95\%。问题二与问题三的精确对照样本分别纳入高负荷、多项目、床旁和专项设备约束；日级完成率与14日滚动率同时保留，用于识别总体率掩盖的短期波动。预测留出、精确子问题、独立可行性复核、单因素敏感性和多种子检验从不同角度限定了模型结果的适用范围。
